\documentclass[10pt,a4paper]{amsart}
\usepackage{amsmath,amssymb,amsthm,graphicx}
\usepackage{hyperref}
\numberwithin{equation}{section}

\newtheorem{Theorem}{Theorem}[section]
\newtheorem*{Theorem*}{Theorem}
\newtheorem{Corollary}[Theorem]{Corollary}
\newtheorem{Proposition}[Theorem]{Proposition}
\newtheorem*{Proposition*}{Proposition}
 { \theoremstyle{definition}
\newtheorem{Example}[Theorem]{Example}
\newtheorem{Remark}[Theorem]{Remark} }

\newcommand{\skal}[2]{\langle #1,#2\rangle}

\title{Minimal G-structures correspond to harmonic sections for the intrinsic-torsion-induced metric}
\author{Kamil Niedzialomski}
\date{doi:10.3842/SIGMA.2016.107; CC BY 4.0}
\begin{document}
\maketitle
\noindent\textit{Allowed standard prerequisites (editorial).}\begin{itemize}
\item Research prerequisites named in the original article. Full original bibliography retained; all source-local definitions and supporting arguments retained.
\end{itemize}\medskip
\section{Introduction}

Existence of a geometric structure on an oriented Riemannian manifold is equivalent to saying that the structure group ${\rm SO}(n)$ of the oriented orthonormal frame bundle reduces to a certain subgroup $G$. For example, for $G$ equal
\begin{gather*}
{\rm SO}(m)\times {\rm SO}(n-m),\quad {\rm U}(n/2),\quad {\rm U}(n/2)\times 1,\quad {\rm Sp}(n/4){\rm Sp}(1),\quad {\rm G}_2,\quad {\rm Spin}(7)
\end{gather*}
we have almost product, almost Hermitian, almost contact, almost quaternion-K\"ahler, ${\rm G}_2$ and ${\rm Spin}(7)$ structures, respectively. It is natural to ask if the holonomy group of the Levi-Civita connection $\nabla$ is contained in given $G$. The list of possible irreducible Riemannian holonomies is limited by the Berger list~\cite{lb}.

On the other hand, the defect of the Levi-Civita connection to be a $G$-connection measures the intrinsic torsion $\xi$, which is the dif\/ference of $\nabla$ and a $G$-connection $\nabla^G$ (with torsion),
\begin{gather*}
\xi_XY=\nabla_XY-\nabla^G_XY.
\end{gather*}
If a $G$-structure is integrable, i.e., the intrinsic torsion vanishes, then the holonomy is contained in the structure group $G$. The study of possible intrinsic torsions, i.e., the decomposition of the space of intrinsic torsions into irreducible modules, was initiated by Gray and Hervella \cite{gh} in the case of almost Hermitian manifolds. Later, many authors considered other possible cases (see, for example, \cite{cg,chs,cs,dm,ffs,fg,mc,ms,an,vw}).

The other possible direction, initiated by Wood \cite{cw3} and generalized to the general case by Gonz{\'a}lez-D{\'a}vila and Mart{\'{\i}}n~Cabrera~\cite{gmc2}, is to consider dif\/ferential properties of intrinsic torsion induced by condition of harmonicity of the unique section of the associated homogeneous bundle. More precisely, a $G$-structure $P\subset {\rm SO}(M)$, where $G$ is closed and connected, induces the unique section $\sigma$ of the homogeneous associated bundle ${\rm SO}(M)/G={\rm SO}(M)\times_{{\rm SO}(n)} {\rm SO}(n)/G$,
\begin{gather*}
\sigma(\pi_{{\rm SO}(M)}(p))=[[p,eG]],\qquad p\in {\rm SO}(M),
\end{gather*}
where $\pi_{{\rm SO}(M)}$ is the projection in the orthonormal frame bundle ${\rm SO}(M)$. Then the decomposition $\mathfrak{so}(n)=\mathfrak{g}\oplus\mathfrak{m}$, where $\mathfrak{m}$ is the orthogonal complement to $\mathfrak{g}$ with respect to the Killing form, on the level of Lie algebras is reductive. Equip ${\rm SO}(M)/G$ with the natural Riemannian metric induced from the Killing form on $\mathfrak{m}$ and Riemannian metric $g$ on $M$. Then we say that a $G$-structure $P$ is harmonic if the section $\sigma$ is harmonic. The correspondence of the notion of harmonicity with the intrinsic torsion follows from the fact that the intrinsic torsion $\xi$ can be considered as a section of the bundle $T^{\ast}M\otimes\mathfrak{m}_P$, where $\mathfrak{m}_P$ is the adjoint bundle $P\times_{\operatorname{ad}G}\mathfrak{m}$. This follows from the observation that the $\mathfrak{m}$-component of the connection form $\omega$ with respect to the reductive decomposition $\mathfrak{so}(n)=\mathfrak{g}\oplus\mathfrak{m}$ can be projected to the tangent bundle $TM$ (see the next section for the details).

To seek for the `best' possible non-integrable $G$-structures we consider the third possible approach. Namely, we focus on the minimality of a $G$-structure in the oriented orthonormal frame bundle ${\rm SO}(M)$. It is not surprising that minimality is related with the harmonicity of a~$G$-structure. More generally, we use the concept of intrinsic torsion to obtain some results on the geometry of~$G$-structures. The idea comes from the results obtained by the author~\cite{kn} in the case of a single submanifold. We deal with the intrinsic and extrinsic geometry of a~$G$-structure. To be more precise, we def\/ine the Riemannian metric on ${\rm SO}(M)$ by inducing it from the Riemannian metric~$g$ on~$M$ and Killing form ${\bf B}$ on the structure group ${\rm SO}(n)$. It is interesting that the Levi-Civita connection on~$P$ depends on the $G$-connection and Levi-Civita connection $\tilde{\nabla}$, which comes from a certain Riemannian metric $\tilde{g}$ on~$M$. This metric $\tilde{g}$ depends on the intrinsic torsion and equals the pull-back of the Riemannian metric on the associated homogeneous bundle ${\rm SO}(M)/G$ with respect to the section $\sigma$.

The main emphasis is put on the minimality of a $G$-structure $P$ in the Riemannian struc\-tu\-re~${\rm SO}(M)$. The main theorem can be stated as follows.
\begin{Theorem*}
The following conditions are equivalent:
\begin{enumerate}\itemsep=0pt
\item[$1)$] $G$-structure $P$ is minimal in ${\rm SO}(M)$,
\item[$2)$] the induced section $\sigma\colon (M,\tilde{g})\to (N,\skal{\cdot}{\cdot})$ is a harmonic map, where $\tilde{g}$ is such that $\tilde{g}=\sigma^{\ast}\skal{\cdot}{\cdot}$.
\end{enumerate}
\end{Theorem*}
Notice, that these conditions are also equivalent, by the general fact concerning harmonic maps \cite{el}, to the minimality of $\sigma(M)$ in $N$. Moreover, the vanishing of the second fundamental form is the neccesary condition for the integrability of a $G$-structure.
\begin{Proposition*}
If a $G$-structure $P$ is integrable, i.e., the intrinsic torsion vanishes, then $P$ is totally geodesic in ${\rm SO}(M)$.
\end{Proposition*}

The article is organized as follows. In Sections~\ref{section2} and~\ref{section3} we recall the notion of the intrinsic torsion and state its main properties. The results in these sections are well known and can be found in the literature, for example, in~\cite{gmc2,cw2}.

In Section~\ref{section4} we introduce a tensor, which transfers the Riemannian metric $g$ to the mentioned above $\tilde{g}$. Its properties are crucial in the main considerations. With the Riemannian metric~$\tilde{g}$ and the horizontal distribution $\mathcal{H}'$ of the $G$-structure $P$ induced by the minimal $G$-connec\-tion~$\nabla '$, the projection $\pi_P\colon P\to M$ becomes the Riemannian submersion.

Section~\ref{section5} is the main section of the general outline and deals with the geometry of $G$-structures. Firstly, we consider intrinsic geometry focusing on the curvatures and, secondly, we consider extrinsic geometry, with the main result concerning minimality of a $G$-structure.

We end the article with some relevant examples. Examples considered here have been considered by other authors \cite{ggv,gd,hl}, so we only list them and state relevant conditions.

Throughout the paper we will use the following notation and identif\/ication: For any associated bundle $E=P\times_G S$ with the f\/iber $S$ and induced by the principal bundle $P(M,G)$ any element in $E$ will be denoted by $[[p,s]]$. Moreover, we have $\Gamma(E)\equiv C^{\infty}(P,S)^G$, where $C^{\infty}(P,S)^G$ is the space of equivariant functions $f\colon P\to S$, $f(pg)=g^{-1}f(p)$. The identif\/ication is the following
\begin{gather*}
\Gamma(E)\ni\sigma\longleftrightarrow f\in C^{\infty}(P,S)^G,\qquad \sigma(\pi_P(p))=[[p,f(p)]].
\end{gather*}

\section{Intrinsic torsion}\label{section2}

In this section we will review the basic facts concerning intrinsic torsion of a $G$-structure \cite{gmc2,cw2}.

Let $(M,g)$ be an oriented Riemannian manifold, ${\rm SO}(M)$ its oriented orthonormal frame bundle. Denote by $\omega$ the connection form induced by the Levi-Civita connection $\nabla$ on $M$. Let~$\mathcal{H}$ and~$\mathcal{V}$ be horizontal and vertical distributions on ${\rm SO}(M)$, respectively.

Assume that the structure group ${\rm SO}(n)$ reduces to a closed and connected subgroup $G$. Then the quotient~${\rm SO}(n)/G$ is a normal homogeneous space, i.e., the subspace $\mathfrak{m}=\mathfrak{g}^{\bot}\subset \mathfrak{so}(n)$ def\/ines~$\operatorname{ad}(G)$ invariant decomposition $\mathfrak{so}(n)=\mathfrak{g}\oplus \mathfrak{m}$, where $\mathfrak{g}$ is the Lie algebra of $G$ and the orthogonal part is taken with respect to the Killing form ${\bf B}$. Denote by $P\subset {\rm SO}(M)$ the reduced subbundle. The $\mathfrak{g}$-component $\omega_{\mathfrak{g}}$ of $\omega$ def\/ines a connection form on $P$. Denote by $\mathcal{H}'$ and $\mathcal{V}'$ the horizontal and vertical distributions on $P$ with respect to $\omega_{\mathfrak{g}}$, respectively.

Moreover, consider the following associated bundles -- the adjoint bundles over $M$: $\mathfrak{so}(n)_P=P\times_{\operatorname{ad}G}\mathfrak{so}(n)$, $\mathfrak{g}_P=P\times_{\operatorname{ad}G}\mathfrak{g}$ and $\mathfrak{m}_P=P\times_{\operatorname{ad}G}\mathfrak{m}$.

Notice that $\mathfrak{so}(n)_P$ can be realized as the bundle $\mathfrak{so}(M)$ of skew-symmetric endomorphisms of the tangent bundle $TM$ via the identif\/ication
\begin{gather*}
[[p,A]]\mapsto p^{-1}\cdot A\cdot p,
\end{gather*}
where we treat element $p\in P_x$ as a linear isomorphism $p\colon \mathbb{R}^n\to T_xM$. Analogously, we have the bundles $\mathfrak{g}(M)$ and $\mathfrak{m}(M)$. Moreover, $\mathfrak{so}(n)_P$ is isomorphic to the bundle $\mathcal{V}_{\rm equiv}$ over $M$ of equivariant vertical vector f\/ields on $P$. Namely, the map
\begin{gather*}
[[p,A]]\mapsto A^{\ast}_p
\end{gather*}
settles this isomorphism, where $A^{\ast}$ denotes the fundamental vertical vector induced by the matrix $A\in\mathfrak{so}(n)$.

Let
\begin{gather*}
\xi(X)=\omega_{\mathfrak{m}}\big(X^{h'}\big),
\end{gather*}
where $\omega_{\mathfrak{m}}$ denotes the $\mathfrak{m}$-component of $\omega$ and $X^{h'}$ is the horizontal lift with respect to $\mathcal{H}'$. Since
\begin{gather*}
\omega_{\mathfrak{m}}\big(X^{h'}_{pg}\big)=\omega_{\mathfrak{m}}\big(R_{g\ast}X^{h'}_p\big)
=\operatorname{ad}\big(g^{-1}\big)\omega_{\mathfrak{m}}\big(X^{h'}_p\big),
\end{gather*}
it follows that $\xi(X)\in \mathfrak{m}_P$. Section $\xi\in\Gamma(T^{\ast}M\otimes\mathfrak{m}_P)$ is called the {\it intrinsic torsion} of a $G$-structure $P$. By above identif\/ications $\xi(X)=\xi_X\in \mathfrak{m}(M)$. Let $\nabla '$ be the connection on $M$ induced by $\omega_{\mathfrak{g}}$ on $P$. We call $\nabla '$ the {\it minimal connection} of a $G$-structure. Intrinsic torsion $\xi$ has the following properties:
\begin{gather}\label{eq:intrinsictorsionproperties}
X^{h'}=X^h+(\xi_X)^{\ast}\qquad\textrm{and}\qquad \xi_XY=\nabla_XY-\nabla'_XY,
\end{gather}
for $X,Y\in \Gamma(TM)$. The decomposition $\mathfrak{so}(M)=\mathfrak{g}(M)\oplus\mathfrak{m}(M)$ implies the following relations between intrinsic torsion and the curvature tensor $R$ of $\nabla$:
\begin{gather}\label{eq:curvaturedecomp}
\begin{split}
&(\nabla_X\xi_Y)_{\mathfrak{g}}=[\xi_X,\xi_Y]_{\mathfrak{g}},\\
&(\nabla_X\xi_Y)_{\mathfrak{m}}=\nabla '_X\xi_Y+[\xi_X,\xi_Y]_{\mathfrak{m}},\\
&R(X,Y)_{\mathfrak{g}}=R'(X,Y)+[\xi_X,\xi_Y]_{\mathfrak{g}},\\
&R(X,Y)_{\mathfrak{m}}=\nabla '_X\xi_Y-\nabla '_Y\xi_X+[\xi_X,\xi_Y]_{\mathfrak{m}}-\xi_{[X,Y]}.
\end{split}
\end{gather}
Notice also that
\begin{gather}\label{eq:Rdecomp2}
R(X,Y)=R'(X,Y)+(\nabla_X\xi)_Y-(\nabla_Y\xi)_X-[\xi_X,\xi_Y],
\end{gather}
where $R'$ is the curvature tensor of $\nabla'$ and
\begin{gather*}
(\nabla_X\xi)_Y=\nabla_X\xi_Y-\xi_{\nabla_XY}.
\end{gather*}

\section[ Harmonic $G$-reductions]{Harmonic $\boldsymbol{G}$-reductions}\label{section3}

Let $(M,g)$ be an oriented Riemannian manifold, $\pi_{{\rm SO}(M)}\colon {\rm SO}(M)\to M$ its orthonormal frame bundle with the connection form $\omega$ inducing Levi-Civita connection $\nabla$ on $M$. Let $G\subset {\rm SO}(n)$, $n=\dim M$, be a subgroup such that ${\rm SO}(n)/G$ is a normal homogeneous space and let $P\subset {\rm SO}(M)$ the reduced subbundle. Let $N={\rm SO}(M)/G={\rm SO}(M)\times_{{\rm SO}(n)} ({\rm SO}(n)/G)$ be the homogeneous bundle associated with ${\rm SO}(M)$. Denote by $\zeta\colon {\rm SO}(M)\to N$ the natural projection. Then $\zeta$ def\/ines the $G$-principal bundle. Clearly $\zeta$ is constant on $P$ and hence there is a bijection between $G$-reductions of ${\rm SO}(M)$ and sections of $N$. Denote by $\sigma\in \Gamma(N)$ the section induced by~$P$. Let $\mathfrak{m}_N$ be the adjoint bundle associated with $\zeta$, i.e., $\mathfrak{m}_N={\rm SO}(M)\times_{\operatorname{ad}G}\mathfrak{m}$.

Let $\mathcal{H}^N$ and $\mathcal{V}^N$ be horizontal and vertical distributions on $N$, respectively, where $\mathcal{H}^N$ is induced by $\omega$. We have
\begin{gather*}
\mathcal{H}^N=\zeta_{\ast}\mathcal{H},\qquad \mathcal{V}^N={\rm SO}(M)\times_{{\rm SO}(n)} T({\rm SO}(n)/G),
\end{gather*}
where ${\rm SO}(n)$ acts on $T({\rm SO}(n)/G)$ by the dif\/ferential of the natural action of ${\rm SO}(n)$ on ${\rm SO}(n)/G$. Since $T({\rm SO}(n)/G)={\rm SO}(n)\times_{\operatorname{ad}G}\mathfrak{m}$ it follows that $\mathcal{V}^N$ is isomorphic to $\mathfrak{m}_N$.

Denote by $\varphi\colon TN\to \mathfrak{m}_N$ the following map
\begin{gather*}
\varphi([[p,[[g,A]]\,]])=[[p,A]],\qquad \varphi\big(X^{h,N}\big)=0,
\end{gather*}
i.e., $\varphi$ settles the described above isomorphism of $\mathcal{V}^N$ onto $\mathfrak{m}_N$ and is zero on the horizontal distribution. The Riemannian metric on $N$ is induced by $g$ and the Killing form on~$\mathfrak{m}$, namely,
\begin{gather}\label{eq:riemmetriconN}
\skal{V}{W}=g(\pi_{N\ast}V,\pi_{N\ast}W)+{\bf B}(\varphi(V),\varphi(W)),\qquad V,W\in TN,
\end{gather}
where ${\bf B}(A,B)=-\operatorname{tr}(AB)$ for $A,B\in\mathfrak{m}$. Decomposition
\begin{gather*}
TN=\mathcal{H}^N\oplus\mathcal{V}^N,
\end{gather*}
def\/ines projections ${\bf h}\colon TN\to\mathcal{H}^N$ and ${\bf v}\colon TN\to\mathcal{H}^N$. It can be shown that~\cite{gmc2}
\begin{gather}\label{eq:verticalsigma}
{\bf v}\sigma_{\ast}(X)=\xi_X\in\mathfrak{m}_N,\qquad X\in TM,
\end{gather}
where $\xi_X$ is the intrinsic torsion of a $G$-structure $P$.

In this section we will derive the formula for the derivative $\nabla\sigma_{\ast}$ and the tension f\/ield of~$\sigma$.

Equip $N$ with the Riemannian metric $\skal{\cdot}{\cdot}$ given by \eqref{eq:riemmetriconN} and~$M$ with a Riemannian metric~$\tilde{g}$ on~$M$, which may vary from~$g$. We will consider $\sigma$ as a map
\begin{gather*}
\sigma\colon \ (M,\tilde g)\to (N,\skal{\cdot}{\cdot}).
\end{gather*}
Assuming $M$ is compact we may def\/ine the energy functional $E(\sigma)$ given by
\begin{gather}\label{eq:energyfunctional}
E(\sigma)=\frac{1}{2}\int_M \|\sigma_{\ast}\|^2\,d{\rm vol}_M,
\end{gather}
where the norm $\|\cdot\|$ is taken with respect to $\tilde g$ and $\skal{\cdot}{\cdot}$, i.e.,
\begin{gather*}
\|\sigma_{\ast}\|^2=\sum_i\skal{\sigma_{\ast}(\tilde{e_i})}{\sigma_{\ast}(\tilde{e_i})},
\end{gather*}
where $(\tilde{e_i})$ is a $\tilde{g}$-orthonormal basis on $M$. In order to study harmonic sections it is convenient to study variations of the functional~\eqref{eq:energyfunctional} in the class of all sections of $N$. We will consider that above functional in the class of all maps from $M$ to $N$, thus we will study harmonicity as a map (not a section). Then, the Euler--Lagrange equation is given by the formula
\begin{gather}\label{eq:ELequation}
\tau_{\tilde{g}}(\sigma)=\operatorname{tr}_{\tilde g}\nabla\sigma_{\ast}
=\sum_i\nabla^{\sigma}_{\tilde{e_i}}\sigma_{\ast}\tilde{e_i}
-\sigma_{\ast}(\tilde{\nabla}_{\tilde{e_i}}\tilde{e_i})=0,
\end{gather}
where $\tilde{\nabla}$ is the Levi-Civita connection of $\tilde{g}$ and $\nabla^{\sigma}$ is the pull-back connection in the bundle
\begin{gather}\label{eq:TNdecomposition}
\sigma^{\ast}TN= \sigma^{\ast}\mathcal{H}^N\oplus\sigma^{\ast}\mathcal{V}^N =TM\oplus\mathfrak{m}_P=TM\oplus\mathfrak{m}(M).
\end{gather}
We call $\tau(\sigma)$ the {\it tension field} of $\sigma$. We say that a $G$-structure is {\it harmonic as a map} if~\eqref{eq:ELequation} holds for the induced section $\sigma$ (here and furthermore we do not require $M$ to be compact). Taking the decomposition of $\tau(\sigma)$ with respect to~\eqref{eq:TNdecomposition}, harmonicity of $\sigma$ is equivalent to vanishing of~${\bf h}\tau(\sigma)$ and~${\bf v}\tau(\sigma)$.

If $\tilde g=g$, then a $G$-structure induced by the section $\sigma$ satisfying ${\bf v}\tau(\sigma)=0$ is called a {\it harmonic} $G$-{\it structure}~\cite{gmc2}.

Denote by $\Pi_{\tilde g}$ and $\Pi_g$ the dif\/ferential $\nabla\sigma_{\ast}$ with respect to $\tilde g$ and $g$, respectively. Moreover, let $S$ be the dif\/ference between $\tilde{\nabla}$ and $\nabla$, i.e., $S(X,Y)=\tilde{\nabla}_XY-\nabla_XY$. Then
\begin{gather*}
\tau_{\tilde g}(\sigma)=\operatorname{tr}_{\tilde g}\Pi_{\tilde g}
=\operatorname{tr}_{\tilde g}(\Pi_g-\sigma_{\ast}S).
\end{gather*}
We have \cite{gmc2, cw2}
\begin{gather}
\varphi(\Pi_g(X,Y))=\frac{1}{2}((\nabla_X\xi)_Y+(\nabla_Y\xi)_X)\in\mathfrak{m}(M), \label{eq:vtau}\\
g(\pi_{N\ast}\Pi_g(X,Y),Z) =\frac{1}{2}(
{\bf B}(\xi_X,R_{\mathfrak{m}}(Y,Z))+{\bf B}(\xi_Y,R_{\mathfrak{m}}(X,Z))), \label{eq:htau}
\end{gather}
where $(\nabla_X\xi)_Y=\nabla_X\xi_Y-\xi_{\nabla_XY}$.

In order to describe explicitly the condition for the harmonicity of $\sigma$ let us introduce certain curvature operator~\cite{kn}. For any $\alpha\in\mathfrak{so}(M)$ let
\begin{gather}\label{eq:defRT}
R_{\alpha}(X)=\sum_i R(e_i,\alpha(e_i))X\in TM,\qquad X\in TM.
\end{gather}
Then $R_{\alpha}\in\mathfrak{so}(M)$ and the following formula holds
\begin{gather}\label{eq:R_Tproperty}
g(R_{\alpha}(X),Y)={\bf B}(\alpha,R(X,Y)),\qquad \alpha\in\mathfrak{so}(M),\qquad X,Y\in TM.
\end{gather}
Indeed,
\begin{gather*}
g(R_{\alpha}(X),Y) =\sum_i g(R(e_i,\alpha(e_i))X,Y)=\sum_i g(\alpha(e_i),R(X,Y)e_i) ={\bf B}(\alpha,R(X,Y)).
\end{gather*}

Now, we can state the formula for the harmonicity of $\sigma$.
\begin{Proposition}\label{prop:harmonicsigma1}
$G$-structure $\sigma\colon (M,\tilde{g})\to (N,\skal{\cdot}{\cdot})$ is a harmonic map if and only if the following conditions hold
\begin{gather*}%\label{eq:harmonicsigma}
\sum_i (\nabla_{\tilde{e_i}}\xi)_{\tilde{e_i}}-\xi_{S(\tilde{e_i},\tilde{e_i})}=0 \qquad\textrm{and}\qquad
\sum_iR_{\xi_{\tilde{e_i}}}(\tilde{e_i})-S(\tilde{e_i},\tilde{e_i})=0,
\end{gather*}
where $(\tilde{e_i})$ is an orthonormal basis for $\tilde{g}$.
\end{Proposition}
\begin{proof}
Follows by \eqref{eq:verticalsigma}, \eqref{eq:vtau}, \eqref{eq:htau} and \eqref{eq:R_Tproperty}.
\end{proof}

\section{Properties of certain transfer tensor}\label{section4}

In this section we will introduce invertible tensor induced by the intrinsic torsion, the Riemannian metric def\/ined by this tensor and state the properties of the Levi-Civita connection of this new metric. Results in this section are generalizations of the results obtained by the author in \cite{kn}.

Adopt the notation from the previous section. For $\alpha\in \mathfrak{m}(M)$ put
\begin{gather*}
\xi\cdot \alpha=-\sum_i {\bf B}(\xi_{e_i},\alpha)e_i\in TM,
\end{gather*}
where we consider the intrinsic torsion as an element of $\mathfrak{m}(M)$. It is easy to show that
\begin{gather}\label{eq:xidotT}
g(\xi\cdot \alpha,X)=-{\bf B}(\alpha,\xi_X),\qquad \alpha\in\mathfrak{m}(M),\qquad X\in TM.
\end{gather}

Let $L$ be the endomorphism of the tangent bundle of the form
\begin{gather*}
L(X)=X-\xi\cdot \xi_X,\qquad X\in TM.
\end{gather*}
In order to derive some properties of $L$ recall the def\/inition of the Riemannian metric on the bundle ${\rm SO}(M)$ induced by the metric $g$ on $M$ and by the Killing form ${\bf B}$ on the structure group~${\rm SO}(n)$. We def\/ine the Riemannian metric $g_{{\rm SO}(M)}$ on~${\rm SO}(M)$ as follows
\begin{gather*}
g_{{\rm SO}(M)}\big(X^h,Y^h\big) =g(X,Y), \qquad
g_{{\rm SO}(M)}\big(X^h,\alpha^{\ast}\big) =0, \qquad
g_{{\rm SO}(M)}\big(\alpha^{\ast},\beta^{\ast}\big) ={\bf B}(\alpha,\beta),
\end{gather*}
where $X,Y\in TM$, $\alpha,\beta\in\mathfrak{so}(M)=\mathfrak{so}(n)_{{\rm SO}(M)}$. Recall that we identify element $\alpha\in \mathfrak{so}(M)$ with the equivariant vertical vector f\/ield denoted by $\alpha^{\ast}$. Then maps $\pi_{{\rm SO}(M)}\colon {\rm SO}(M)\to M$ and $\zeta\colon {\rm SO}(M)\to N$ are Riemannian submersions.

\begin{Proposition}\label{prop:PropertiesofL}
We have
\begin{gather*}
g_{{\rm SO}(M)}\big(X^{h'},Y^{h'}\big)=g(X,LY).
\end{gather*}
In particular, $L$ is a symmetric and positive definite automorphism of $TM$. Moreover, the covariant derivative of $L$ is related with the intrinsic torsion $\xi_X$ by the formula
\begin{gather*}
g((\nabla_XL)Y,Z)={\bf B}((\nabla_X\xi)_Y,\xi_Z)+{\bf B}((\nabla_X\xi)_Z,\xi_Y).
\end{gather*}
\end{Proposition}
\begin{proof}
By \eqref{eq:xidotT} and \eqref{eq:intrinsictorsionproperties},
\begin{gather*}
g(X,LY) =g(X,Y)+{\bf B}(\xi_X,\xi_Y)=g_{{\rm SO}(M)}\big(X^h+\xi_X,Y^h+\xi_Y\big) =g_{{\rm SO}(M)}\big(X^{h'},Y^{h'}\big).
\end{gather*}
Thus $L$ is symmetric and positive def\/inite. Moreover, by the fact that $\nabla$ is metric for ${\bf B}$ we have
\begin{gather*}
g((\nabla_XL)Y,Z)=g(\nabla_X(LY),Z)-g(L(\nabla_XY),Z)\\
\hphantom{g((\nabla_XL)Y,Z)}{} =Xg(LY,Z)-g(LY,\nabla_XZ)-g(\nabla_XY,LZ)\\
\hphantom{g((\nabla_XL)Y,Z)}{}=Xg(Y,Z)+X{\bf B}(\xi_Y,\xi_Z)-g(Y,\nabla_XZ)-{\bf B}(\xi_Y,\xi_{\nabla_XZ})\\
\hphantom{g((\nabla_XL)Y,Z)=}{}-g(\nabla_XY,Z)-{\bf B}(\xi_{\nabla_XY},\xi_Z)\\
\hphantom{g((\nabla_XL)Y,Z)}{}={\bf B}(\nabla_X\xi_Y-\xi_{\nabla_XY},\xi_Z) +{\bf B}(\nabla_X\xi_Z-\xi_{\nabla_XZ},\xi_Y).\tag*{\qed}
\end{gather*}\renewcommand{\qed}{}
\end{proof}

By Proposition \ref{prop:PropertiesofL} the symmetric and bilinear form
\begin{gather*}
\tilde{g}(X,Y)=g(X,LY),\qquad X,Y\in TM,
\end{gather*}
def\/ines a Riemannian metric on $M$. We call $L$ the {\it transfer tensor} between $g$ and~$\tilde g$. Notice, that the projection $\pi_P\colon P\to M$ is a Riemannian submersion with respect to $g_{{\rm SO}(M)}$ on~$P$ and~$\tilde{g}$ on~$M$. Denote by $\tilde{\nabla}$ the Levi-Civita connection of~$\tilde{g}$. One can show that \cite{gm}
\begin{gather*}%\label{eq:differenceofnablas}
2g\big(\tilde{\nabla}_XY-\nabla_XY,LZ\big)=g((\nabla_XL)Y,Z)+g((\nabla_YL)X,Z)-g(X,(\nabla_ZL)Y).
\end{gather*}
Thus by Proposition \ref{prop:PropertiesofL} we get
\begin{gather}
2g\big(\tilde{\nabla}_XY-\nabla_XY,LZ\big) ={\bf B}((\nabla_X\xi)_Y+(\nabla_Y\xi)_X,\xi_Z)\nonumber\\
\hphantom{2g\big(\tilde{\nabla}_XY-\nabla_XY,LZ\big) =}{} +{\bf B}((\nabla_X\xi)_Z-(\nabla_Z\xi)_X,\xi_Y)+{\bf B}((\nabla_Y\xi)_Z-(\nabla_Z\xi)_Y,\xi_X).\label{eq:differencetensor}
\end{gather}

\section[Geometry of $G$-structures]{Geometry of $\boldsymbol{G}$-structures}\label{section5}

In this section we study the geometry of a $G$-structure $P$ in ${\rm SO}(M)$. In this case we need to consider the Levi-Civita connection $\nabla^{{\rm SO}(M)}$ of the Riemannian metric~$g_{{\rm SO}(M)}$. One can show~\cite{kn} that
\begin{alignat*}{3}
& \nabla^{{\rm SO}(M)}_{X^h}X^h =(\nabla_XY)^h-\frac{1}{2}R(X,Y)^{\ast},\qquad && \nabla^{{\rm SO}(M)}_{\alpha^{\ast}}Y^h =\frac{1}{2}R_{\alpha}(Y)^h,&\\
& \nabla^{{\rm SO}(M)}_{X^h}\alpha^{\ast} =\frac{1}{2}R_{\alpha}(X)^h+(\nabla_X\alpha)^{\ast},\qquad &&
\nabla^{{\rm SO}(M)}_{\alpha^{\ast}}\beta^{\ast} =-\frac{1}{2}[\alpha,\beta]^{\ast},&
\end{alignat*}
for $X,Y\in \Gamma(TM)$ and $\alpha,\beta\in\mathfrak{so}(M)$, or equivalently, $\alpha,\beta\in \mathfrak{so}(n)_{{\rm SO}(M)}$, where $R_{\alpha}$ is def\/ined by~\eqref{eq:defRT}.

It is convenient to f\/ind orthogonal projections
\begin{gather*}
T{\rm SO}(M)\mapsto TPq\quad\textrm{and}\qquad T{\rm SO}(M)\mapsto T^{\bot}P,
\end{gather*}
in order to derive the Levi-Civita connection and second fundamental form for $P$ in ${\rm SO}(M)$.

\begin{Proposition}\label{prop:ProjtoTP}
Let $X\in TM$ and $\alpha\in\mathfrak{so}(M)$.
\begin{enumerate}\itemsep=0pt
\item[$1.$] The orthogonal projection $T{\rm SO}(M)\mapsto TP$ equals
\begin{gather*}
X^h\mapsto L^{-1}(X)^{h'},\qquad
\alpha^{\ast}\mapsto \alpha_{\mathfrak{g}}^{\ast}-L^{-1}(\xi\cdot\alpha_{\mathfrak{m}})^{h'}.
\end{gather*}
\item[$2.$] The orthogonal projection $T{\rm SO}(M)\mapsto T^{\bot}P$ equals
\begin{gather*}
X^h\mapsto -(\xi_{L^{-1}(X)})^{\ast}-(\xi\cdot\xi_{L^{-1}X})^h,
\qquad
\alpha^{\ast}\mapsto \alpha_{\mathfrak{m}}^{\ast}
+(\xi_{L^{-1}(\xi\cdot\alpha_{\mathfrak{m}})})^{\ast}
+L^{-1}(\xi\cdot\alpha_{\mathfrak{m}})^h.
\end{gather*}
\end{enumerate}
\end{Proposition}
\begin{proof}
We will use frequently \eqref{eq:intrinsictorsionproperties} and \eqref{eq:xidotT}. Recall that $TP=\mathcal{H}'\oplus\mathfrak{m}_P$. For any $Y\in TM$
\begin{gather*}
g_{{\rm SO}(M)}\big(X^h-L^{-1}(X)^{h'},Y^{h'}\big) =g_{{\rm SO}(M)}\big(X^h-L^{-1}(X)^h-\xi_{L^{-1}(X)}^{\ast},Y^h+\xi_Y\big)\\
\hphantom{g_{{\rm SO}(M)}\big(X^h-L^{-1}(X)^{h'},Y^{h'}\big)}{} =g\big(X-L^{-1}(X),Y\big)-{\bf B}\big(\xi_{L^{-1}(X)},\xi_Y\big)=0
\end{gather*}
and for any $\beta\in\mathfrak{g}_P$
\begin{gather*}
g_{{\rm SO}(M)}\big(X^h-L^{-1}(X)^{h'},\beta^{\ast}\big)=-{\bf B}\big(\xi_{L^{-1}(X)},\beta\big)=0.
\end{gather*}
Clearly $L^{-1}(X)$ is tangent to $P$ and $X^h-L^{-1}(X)^{h'}=-(\xi_{L^{-1}(X)})^{\ast}-(\xi\cdot\xi_{L^{-1}X})^h$, which proves the desired decomposition for $X^h$. Analogously, \begin{gather*}
\alpha^{\ast}-\big(\alpha_{\mathfrak{g}}^{\ast}-L^{-1}(\xi\cdot\alpha_{\mathfrak{m}})^{h'}\big)=
\alpha_{\mathfrak{m}}^{\ast}+L^{-1}(\xi\cdot\alpha_{\mathfrak{m}})^{h'},
\end{gather*}
and
\begin{gather*}
g_{{\rm SO}(M)}\big(\alpha_{\mathfrak{m}}^{\ast}+L^{-1}(\xi\cdot\alpha_{\mathfrak{m}})^{h'},Y^{h'}\big) =
{\bf B}(\alpha_{\mathfrak{m}},\xi_Y)+g\big(L^{-1}(\xi\cdot\alpha_{\mathfrak{m}}),Y\big)+{\bf B}\big(\xi_{L^{-1}(\xi\cdot\alpha_{\mathfrak{m}})},\xi_Y\big)\\
\hphantom{g_{{\rm SO}(M)}\big(\alpha_{\mathfrak{m}}^{\ast}+L^{-1}(\xi\cdot\alpha_{\mathfrak{m}})^{h'},Y^{h'}\big)}{}
=-g(\xi\cdot\alpha_{\mathfrak{m}},Y)
+g\big(L^{-1}(\xi\cdot\alpha_{\mathfrak{m}},Y)\big)\!-g\big(\xi\cdot\xi_{L^{-1}(\xi\cdot\alpha_{\mathfrak{m}})},Y\big)\!\\
\hphantom{g_{{\rm SO}(M)}\big(\alpha_{\mathfrak{m}}^{\ast}+L^{-1}(\xi\cdot\alpha_{\mathfrak{m}})^{h'},Y^{h'}\big)}{}=0.
\end{gather*}
Moreover, for $\beta\in\mathfrak{g}_P$
\begin{gather*}
g_{{\rm SO}(M)}\big(\alpha_{\mathfrak{m}}^{\ast}+L^{-1}(\xi\cdot\alpha_{\mathfrak{m}})^{h'},\beta^{\ast}\big)=
{\bf B}(\alpha_{\mathfrak{m}},\beta)
+{\bf B}\big(\xi_{L^{-1}(\xi\cdot\alpha_{\mathfrak{m}})},\beta\big)=0.
\end{gather*}
Proposition follows from the fact that
\begin{gather*}
L^{-1}(\xi\cdot\alpha_{\mathfrak{m}})^{h'}=L^{-1}(\xi\cdot\alpha_{\mathfrak{m}})^h+\big(\xi_{L^{-1}(\xi\cdot\alpha_{\mathfrak{m}})}\big)^{\ast}
\end{gather*}
and $\alpha_{\mathfrak{g}}^{\ast}-L^{-1}(\xi\cdot\alpha_{\mathfrak{m}})^{h'}$ is tangent to $P$.
\end{proof}

\subsection{Intrinsic geometry}

For the intrinsic geometry we will compute the Levi-Civita connection $\nabla^P$ of $(P,g_{{\rm SO}(M)})$, the curvature tensor, Ricci tensor and sectional and scalar curvatures. We will compare geometry of $P$ with the geometry of the base manifold~$M$.

Let us f\/irst introduce operator $Q_{\alpha}$ and establish some relations. For $X\in TM$ and $\alpha\in\mathfrak{so}(M)$ put
\begin{gather*}
Q_{\alpha}(X)=L^{-1}(R_{\alpha}(X)-\xi\cdot(\nabla_X\alpha)_{\mathfrak{m}}).
\end{gather*}
Moreover, it is easy to see that for any $\alpha\in\mathfrak{so}(M)$ we have the following relations with respect to the decomposition $\mathfrak{so}(M)=\mathfrak{g}(M)\oplus\mathfrak{m}(M)$:
\begin{gather}\label{eq:basicrelations}
\begin{split}
&\nabla_X\alpha_{\mathfrak{g}} =\nabla'_X\alpha_{\mathfrak{g}}+[\xi_X,\alpha_{\mathfrak{g}}],\\
&\nabla_X\alpha_{\mathfrak{m}} =[\xi_X,\alpha_{\mathfrak{m}}]_{\mathfrak{g}}+\big( \nabla'_X\alpha_{\mathfrak{m}}+[\xi_X,\alpha_{\mathfrak{m}}]_{\mathfrak{m}} \big).
\end{split}
\end{gather}

\begin{Theorem}\label{thm:nablaP}
The Levi-Civita connection $\nabla^P$ of $P$ with the induced Riemannian metric $g_{{\rm SO}(M)}$ from ${\rm SO}(M)$ is of the following form
\begin{alignat*}{3}
& \nabla^P_{X^{h'}}Y^{h'} =(\tilde\nabla_XY)^{h'}-\frac{1}{2}R'(X,Y)^{\ast},\qquad && \nabla^P_{\alpha^{\ast}}Y^{h'} =\frac{1}{2}Q_{\alpha}(Y)^{h'},&\\
& \nabla^P_{X^{h'}}\beta^{\ast} =\frac{1}{2}Q_{\beta}(X)^{h'}+(\nabla'_X\beta)^{\ast},\qquad && \nabla^P_{\alpha^{\ast}}\beta^{\ast} =-\frac{1}{2}[\alpha,\beta]^{\ast},&
\end{alignat*}
where $X,Y\in TM$, $\alpha,\beta\in\mathfrak{g}_P$.
\end{Theorem}
\begin{proof}
First, by \eqref{eq:xidotT} and \eqref{eq:R_Tproperty} we have
\begin{gather*}
\tilde{g}(Q_{\alpha}(X),Z)={\bf B}(R(X,Z),\alpha)+{\bf B}(\nabla_X\alpha,\xi_Y).
\end{gather*}
Thus, using \eqref{eq:Rdecomp2} and \eqref{eq:differencetensor}
\begin{gather*}
\tilde{g}(Q_{\xi_X}(Y)+Q_{\xi_Y}(X)+\xi\cdot\xi_{\nabla_XY+\nabla_YX},Z)\\
\qquad {}= {\bf B}((\nabla_X\xi)_Z-(\nabla_Z\xi)_X,\xi_Y)+{\bf B}((\nabla_Y\xi)_Z-(\nabla_Z\xi)_Y,\xi_X)-{\bf B}([\xi_X,\xi_Z],\xi_Y)\\
\qquad \quad{}-{\bf B}([\xi_Y,\xi_Z],\xi_X)+{\bf B}((\nabla_Y\xi)_X+(\nabla_X\xi)_Y,\xi_Z)\\
\qquad {} =2\tilde{g}(\tilde{\nabla}_XY-\nabla_XY,Z).
\end{gather*}
We have shown that
\begin{gather}\label{eq:nablaPproof}
P(X,Y)=\frac{1}{2}\big(Q_{\xi_X}(Y)+Q_{\xi_Y}(X)+\xi\cdot\xi_{\nabla_XY+\nabla_YX}\big).
\end{gather}
By the formula for the connection $\nabla^{{\rm SO}(M)}$
\begin{gather*}
\nabla^P_{X^{h'}}Y^{h'} = \nabla^{{\rm SO}(M)}_{X^h}Y^h+\nabla^{{\rm SO}(M)}_{(\xi_X)^{\ast}}Y^h
+\nabla^{{\rm SO}(M)}_{X^h}(\xi_Y)^{\ast}+\nabla^{{\rm SO}(M)}_{(\xi_X)^{\ast}}(\xi_Y)^{\ast}\\
\hphantom{\nabla^P_{X^{h'}}Y^{h'}}{}
=\left(\nabla_XY+\frac{1}{2}R_{\xi_Y}(X)+\frac{1}{2}R_{\xi_X}(Y)\right)^h+\left(-\frac{1}{2}R(X,Y)+\nabla_X\xi_Y-\frac{1}{2}[\xi_X,\xi_Y]\right)^{\ast}.
\end{gather*}
Since, by relations \eqref{eq:curvaturedecomp} between intrinsic torsion and curvature tensor,
\begin{gather*}
\left(-\frac{1}{2}R(X,Y)+\nabla_X\xi_Y-\frac{1}{2}[\xi_X,\xi_Y]\right)_{\mathfrak{g}}=-\frac{1}{2}R'(X,Y)
\end{gather*}
and
\begin{gather*}
\left(-\frac{1}{2}R(X,Y)+\nabla_X\xi_Y-\frac{1}{2}[\xi_X,\xi_Y]\right)_{\mathfrak{m}}=\frac{1}{2}\left( \nabla_X\xi_Y+\nabla_Y\xi_X+\xi_{[X,Y]} \right),
\end{gather*}
then, with the use of Proposition~\ref{prop:ProjtoTP}, we get
\begin{gather*}
\nabla^P_{X^{h'}}Y^{h'}=L^{-1}\left(\frac{1}{2}R_{\xi_Y}(X)+\frac{1}{2}R_{\xi_X}(Y)-\frac{1}{2}\xi\cdot(\nabla_X\xi_Y+\nabla_Y\xi_X)\right)^{h'}\\
\hphantom{\nabla^P_{X^{h'}}Y^{h'}=}{}
+L^{-1}\left(\nabla_XY-\frac{1}{2}\xi\cdot\xi_{[X,Y]}\right)^{h'}-\frac{1}{2}R'(X,Y)^{\ast}.
\end{gather*}
Moreover $\nabla_XY-\frac{1}{2}\xi\cdot\xi_{[X,Y]}=L(\nabla_XY)+\xi\cdot\xi_{\nabla_XY+\nabla_YX}$. Thus, by the def\/inition of $Q_{\alpha}$ and formula~\eqref{eq:nablaPproof} we get the desired formula for $\nabla^P_{X^{h'}}Y^{h'}$.

For the proof of the second formula we will use \eqref{eq:basicrelations} and Proposition~\ref{prop:ProjtoTP}. We have
\begin{gather*}
\nabla^P_{X^{h'}}\beta^{\ast} =\left( \nabla^{{\rm SO}(M)}_{X^h}\beta^{\ast}
+\nabla^{{\rm SO}(M)}_{(\xi_X)^{\ast}}\beta^{\ast} \right)^{\top} =\left(\frac{1}{2}R_{\beta}(X)^h+(\nabla_X\beta)^{\ast}-\frac{1}{2}[\xi_X,\beta]^{\ast}\right)^{\top}\\
\hphantom{\nabla^P_{X^{h'}}\beta^{\ast}}{}
=L^{-1}\left(\frac{1}{2}R_{\beta}(X)-\xi\cdot(\nabla_X\beta)_{\mathfrak{m}}+\frac{1}{2}\xi\cdot[\xi_X,\beta]_{\mathfrak{m}}\right)^{h'}
+(\nabla_X\beta)^{\ast}_{\mathfrak{g}}\\
\hphantom{\nabla^P_{X^{h'}}\beta^{\ast}}{}
=Q_{\beta}(X)^{h'}+(\nabla'_X\beta)^{\ast}.
\end{gather*}
We prove the remaining relations analogously.
\end{proof}

\begin{Proposition}\label{prop:RP}
The curvature tensor $R^P$ of $\nabla^P$ equals
\begin{gather*}
R^P\big(X^{h'},Y^{h'}\big)Z^{h'} =\big(\tilde{R}(X,Y)Z\big)^{h'}-\frac{1}{4}\big( Q_{R'(Y,Z)}(X)-Q_{R'(X,Z)}(Y)-2Q_{R(X,Y)}(Z)\big)^{h'}\\
\hphantom{R^P\big(X^{h'},Y^{h'}\big)Z^{h'} =}{}-\frac{1}{2}(D_XR')(Y,Z)^{\ast}+\frac{1}{2}(D_YR')(X,Z)^{\ast},\\
R^P\big(X^{h'},Y^{h'}\big)\gamma^{\ast} =\frac{1}{2}\big((D_XQ)_{\gamma}(Y)-(D_YQ)_{\gamma}(X)\big)^{h'}+\frac{1}{2}[R'(X,Y),\gamma]^{\ast}\\
\hphantom{R^P\big(X^{h'},Y^{h'}\big)\gamma^{\ast} =}{}
-\frac{1}{4}\big( R'(X,Q_{\gamma}(Y))-R'(Y,Q_{\gamma}(X)) \big)^{\ast},\\
R^P\big(X^{h'},\beta^{\ast}\big)Z^{h'} =\frac{1}{2}(D_XQ)_{\beta}(Z)^{h'}-\frac{1}{4}\big( R'(X,Q_{\beta}(Z))+[\beta,R'(X,Z)] \big)^{\ast},\\
R^P\big(X^{h'},\beta^{\ast}\big)\gamma^{\ast} =-\frac{1}{4}\big( Q_{[\beta,\gamma]}(X)+Q_{\beta}(Q_{\gamma}(X)) \big)^{h'},\\
R^P(\alpha^{\ast},\beta^{\ast})Z^{h'} =\frac{1}{4}[Q_{\alpha},Q_{\beta}](Z)^{h'}+\frac{1}{2}Q_{[\alpha,\beta]}(Z)^{h'},\\
R^P(\alpha^{\ast},\beta^{\ast})\gamma^{\ast} =-\frac{1}{4}[[\alpha,\beta],\gamma]^{\ast},
\end{gather*}
where
\begin{gather*}
(D_XR')(Y,Z) =\nabla '_XR'(Y,Z)-R'(\tilde{\nabla}_XY,Z)-R'(Y,\tilde{\nabla}_XZ),\\
(D_XQ)_{\alpha}(Y) =\tilde{\nabla}_XQ_{\alpha}(Y)-Q_{\nabla '_X\alpha}(Y)-Q(\tilde{\nabla}_XY).
\end{gather*}
\end{Proposition}
\begin{proof}
Follows directly by Theorem \ref{thm:nablaP}, Jacobi identity and the relations
\begin{gather*}
R'(X,Y)\gamma=[R'(X,Y),\gamma],\qquad \nabla '_X[\beta,\gamma]=[\nabla '_X\beta,\gamma]+[\beta,\nabla '_X\gamma].
\end{gather*}
Notice that $[X^{h'},\beta^{\ast}]=(\nabla '_X\beta)^{\ast}$ and $[\alpha^{\ast},\beta^{\ast}]=-[\alpha,\beta]^{\ast}$.
\end{proof}

\begin{Corollary}\label{cor:RicciP}
The Ricci curvature tensor ${\rm Ric}^P$ equals
\begin{gather*}
{\rm Ric}^P\big(X^{h'},Y^{h'}\big) =\widetilde{{\rm Ric}}(X,Y)-\frac{3}{4}\sum_i{\bf B}\big(R'(X,\tilde{e_i}),R'(Y,\tilde{e_i})\big)
+\frac{1}{4}\sum_A \tilde{g}(Q_{\alpha_A}(X),Q_{\alpha_A}(Y)),\\
{\rm Ric}^P(X^{h'},\gamma^{\ast}) =\frac{1}{2}\big( \big(\widetilde{\operatorname{div}}\, Q\big)_{\gamma}(X)-\operatorname{tr}_{\tilde{g}}(D_XQ)_{\gamma} \big),\\
{\rm Ric}^P(\beta^{\ast},\gamma^{\ast}) =\frac{1}{4}\left( \sum_i \tilde{g}\big(Q_{\beta}(\tilde{e_i}),Q_{\gamma}(\tilde{e_i})\big)+\sum_A {\bf B}([\alpha_A,\beta],[\alpha_A,\gamma]) \right),
\end{gather*}
where
\begin{gather*}
\big(\widetilde{\operatorname{div}}\,  Q\big)_{\gamma}(X)=\sum_i\tilde{g}\big((D_{\tilde e_i}Q)_{\gamma}(X),\tilde{e_i}\big)
\end{gather*}
and $(\tilde{e_i})$ is an orthonormal basis with respect to $\tilde{g}$ on $M$ and $(\alpha_A)$ is an orthonormal basis of $\mathfrak{g}_P$ with respect to ${\bf B}$.
\end{Corollary}
\begin{proof}
First, notice that
\begin{gather}\label{eq:QalphaRprime}
\tilde{g}(Q_{\alpha}(X),Y)={\bf B}(R'(X,Y),\alpha),
\end{gather}
which follows by the def\/inition of $Q_{\alpha}$, \eqref{eq:basicrelations} and relations between intrinsic torsion and curvature tensor \eqref{eq:curvaturedecomp}. Hence $Q_{\alpha}$ is skew-symmetric with respect to $\tilde{g}$. Now, it suf\/f\/ices to use the formulas for the curvature tensor~$R^P$.
\end{proof}

\begin{Corollary}\label{cor:SectionalP}
The sectional curvatures of $\nabla^P$ are given by the following formulas
\begin{gather*}
\kappa^P\big(X^{h'},Y^{h'}\big) =\tilde{\kappa}(X,Y)-\frac{3}{4}\|R'(X,Y)\|^2,\\
\kappa^P\big(X^{h'},\beta^{\ast}\big) =\frac{1}{4}\|Q_{\beta}(X)\|^2_{\tilde{g}},\qquad
\kappa^P(\alpha^{\ast},\beta^{\ast}) =\frac{1}{4}\|[\alpha,\beta]\|^2,
\end{gather*}
where $X$, $Y$ are orthonormal with respect to $\tilde{g}$ and $\alpha,\beta\in\mathfrak{g}_P$ orthonormal with respect to~${\bf B}$.
\end{Corollary}
\begin{proof}
First and last relations follow immediately by the formulas for the curvature tensor $R^P$ and by \eqref{eq:QalphaRprime}. For the proof of the second one it suf\/f\/ices to use the skew-symmetry of $Q_{\beta}$ with respect to $\tilde{g}$ (see the proof of Corollary \ref{cor:RicciP}).
\end{proof}

\begin{Corollary}\label{cor:scalarP}
The scalar curvature of $\nabla^P$ equals
\begin{gather*}
s^P=\tilde{s}-\frac{3}{4}\sum_{i,j}\|R'(\tilde{e_i},\tilde{e_j})\|^2+\frac{1}{2}\sum_{i,a}\|Q_{\alpha_A}(\tilde{e_i})\|^2_{\tilde{g}}
+\frac{1}{4}\sum_{A,B}\|[\alpha_A,\alpha_B]\|^2.
\end{gather*}
\end{Corollary}

By the above formulas for the curvatures we have the following relations between the geometry of $(P,g_{{\rm SO}(M)})$ and $M$.

\begin{Corollary}\label{cor:relPandtildeM}
We have:
\begin{enumerate}\itemsep=0pt
\item[$1.$] If $\dim M>2$, then $P$ is never of constant sectional curvature.
\item[$2.$]  If the intrinsic torsion $\xi$ vanishes and $(M,g)$ is of constant sectional curvature $0\leq \kappa\leq\frac{2}{3}$, then $P$ has non-negative sectional curvatures.
\item[$3.$]  ${\rm Ric}^P(\alpha^{\ast})$ is non-negative in any direction $\alpha\!\in\! \mathfrak{g}_P$. Moreover, if $\widetilde{\rm Ric}(X)\!\geq\! \frac{3}{4}\!\sum_i\|R'(X,\tilde{e_i})\|^2$, then ${\rm Ric}^P(X)$ is non-negative.
\item[$4.$]  If $R'=0$ and scalar curvature of $(M,\tilde{g})$ is positive, then the scalar curvature of $P$ is positive.
\end{enumerate}
\end{Corollary}

\subsection{Extrinsic geometry}

The properties of the extrinsic geometry are encoded in the second fundamental form, which we will derive explicitly. Moreover, we will compute the mean curvature vector of $P$ in ${\rm SO}(M)$ and relate the minimality of $P$ with the harmonicity of an induced section $\sigma$ with appropriate Riemannian structure.

Adopt the notation from the previous sections. For $\alpha\in\mathfrak{m}_P$ let
\begin{gather*}
\alpha^{+}=\alpha^{\ast}+(\xi\cdot\alpha)^h.
\end{gather*}
Then $\alpha^{+}\in T^{\bot}P$.

\begin{Theorem}\label{thm:secfundformP}
The second fundamental form $\Pi^P$ of $P$ in ${\rm SO}(M)$ satisfies the following relations
\begin{gather*}
g_{{\rm SO}(M)}\big(\Pi^P(X^{h'},Y^{h'}),\alpha^{+}\big) =\frac{1}{2}{\bf B}((\nabla_X\xi)_Y+(\nabla_Y\xi)_X-\xi_{R_{\xi_X}(Y)+R_{\xi_Y}(X)},\alpha),\\
g_{{\rm SO}(M)}\big(\Pi^P(X^{h'},\gamma^{\ast}),\alpha^{+}\big) =\frac{1}{2}{\bf B}([\xi_X,\gamma]_{\mathfrak{m}}-\xi_{R_{\gamma}(X)},\alpha),\\
g_{{\rm SO}(M)}\big(\Pi^P(\beta^{\ast},\gamma^{\ast}),\alpha^{+}\big)=0.
\end{gather*}
\end{Theorem}
\begin{proof}
By the relations \eqref{eq:curvaturedecomp} (see also proof of Theorem \ref{thm:nablaP}) we have
\begin{gather*}
\nabla^{{\rm SO}(M)}_{X^h}Y^h  =\left( \nabla_XY+\frac{1}{2}R_{\xi_X}(Y)+\frac{1}{2}R_{\xi_Y}(X) \right)^h\\
\hphantom{\nabla^{{\rm SO}(M)}_{X^h}Y^h  =}{} +\frac{1}{2}\big( \nabla_X\xi_Y+\nabla_Y\xi_X+\xi_{[X,Y]} \big)^{\ast}+\textrm{$\mathfrak{g}_P$-component}.
\end{gather*}
Thus
\begin{gather*}
g_{{\rm SO}(M)}\big(\Pi^P(X^{h'},Y^{h'}),\alpha^{+}\big) =g\left(\nabla_XY+\frac{1}{2}(R_{\xi_X}(Y)+R_{\xi_Y}(X)),\xi\cdot\alpha\right)\\
\hphantom{g_{{\rm SO}(M)}\big(\Pi^P(X^{h'},Y^{h'}),\alpha^{+}\big) =}{}
+\frac{1}{2}{\bf B}\big(\nabla_X\xi_Y+\nabla_Y\xi_X+\xi_{[X,Y]},\alpha\big),
\end{gather*}
which implies the f\/irst equality. Moreover,
\begin{gather*}
\nabla^{{\rm SO}(M)}_{X^{h'}}\gamma^{\ast}=\frac{1}{2}R_{\gamma}(X)^h+
\frac{1}{2}[\xi_X,\gamma]^{\ast}+\textrm{$\mathfrak{g}_P$-component},
\end{gather*}
which proves the second equality. Since $[\mathfrak{g},\mathfrak{g}]\subset\mathfrak{g}$, it follows that $\mathfrak{m}_P$-component of $\nabla^{{\rm SO}(M)}_{\alpha^{\ast}}\beta^{\ast}$ vanishes.
\end{proof}

By above theorem we get the following implication.

\begin{Corollary}\label{cor:totgeodP}
If a $G$-structure $P$ is integrable, i.e., the intrinsic torsion $\xi$ vanishes, then $P$ is totally geodesic in ${\rm SO}(M)$.
\end{Corollary}

Now, we can state the main theorem of this article.


\begin{Theorem}\label{thm:minimalityofP}
A $G$-structure $P$ is minimal in ${\rm SO}(M)$ if and only the induced section \mbox{$\sigma\colon M{\to} N$} is a harmonic map with respect to Riemannian metrics $\tilde{g}$ and $\skal{\cdot}{\cdot}$ on $M$ and $N$, respectively.
\end{Theorem}

\begin{proof}
Recall, that by Proposition~\ref{prop:harmonicsigma1} section $\sigma\colon (M,\tilde{g})\to (N,\skal{\cdot}{\cdot})$ is a harmonic map if and only if
\begin{gather}
\sum_i(\nabla_{\tilde{e_i}}\xi)_{\tilde{e_i}}-\xi_{S(\tilde{e_i},\tilde{e_i})}=0,\label{eq:H1}\\ %\tag{H1}\\
\sum_i R_{\xi_{\tilde{e_i}}}(\tilde{e_i})-S(\tilde{e_i},\tilde{e_i})=0,\label{eq:H2} %\tag{H2}
\end{gather}
whereas, by Theorem \ref{thm:secfundformP}, minimality of $P$ is equivalent to the following condition
\begin{gather}\label{eq:minimalityP}%\tag{M}
\sum_i (\nabla_{\tilde{e_i}}\xi)_{\tilde{e_i}}-\xi_{R_{\xi_{\tilde{e_i}}}(\tilde{e_i})}=0.
\end{gather}
Clearly, \eqref{eq:H1} and \eqref{eq:H2} imply \eqref{eq:minimalityP}. Conversely, assume \eqref{eq:minimalityP} holds. It suf\/f\/ices to show that~\eqref{eq:H2} holds. By~\eqref{eq:differencetensor} and~\eqref{eq:minimalityP} we have
\begin{gather*}
g\left( \sum_i R_{\xi_{\tilde{e_i}}}(\tilde{e_i})-S(\tilde{e_i},\tilde{e_i}),LZ\right)\\
\qquad{}=-\sum_i{\bf B}\big((\nabla_{e_i}\xi)_{e_i},\xi_Z\big)
+{\bf B}\big((\nabla_{e_i}\xi)_Z-(\nabla_Z\xi)_{e_i},\xi_{e_i}\big) +\sum_i g\big(R_{\xi_{e_i}}(e_i),Z-\xi\cdot \xi_Z\big)\\
\qquad{} =-\sum_i{\bf B}\big((\nabla_{e_i}\xi)_Z-(\nabla_Z\xi)_{e_i},\xi_{e_i}\big)+{\bf B}\big(R(e_i,Z),\xi_{e_i}\big).
\end{gather*}
By the relations \eqref{eq:curvaturedecomp} and the fact that ${\rm SO}(n)/G$ is a normal homogeneous space we get the desired equality.
\end{proof}


\begin{Remark}\label{rem:H1impliesH2}
Let us comment on Theorem \ref{thm:minimalityofP}.
\begin{enumerate}\itemsep=0pt
\item Two equivalent conditions of Theorem \ref{thm:minimalityofP} are also equivalent to minimality the ima\-ge~$\sigma(M)$ inside $N$. This follows by the general facts concerning harmonic maps~\cite{el}.
\item As it was pointed out by an anonymous referee, condition~\eqref{eq:H1} implies~\eqref{eq:H2} (for $\tilde{g}=\sigma^{\ast}\skal{\,}{\,}$). The proof follows by similar arguments as in the proof of \cite[Proposition~2.5]{ggv}. The author wishes to thank anonymous referee for this observation.
\end{enumerate}
\end{Remark}

\section{Some examples concerning almost product structures}\label{section6}

In this section we illustrate obtained results for the ${\rm SO}(m)\times {\rm SO}(n-m)$-structures, often called almost product structures. The case of other possible $G$-structures, for example coming from the Berger list of possible Riemannian holonomy groups, will be studied by the author independently. The considered examples have been already studied in a similar context \cite{ggv,gd,hl}. Thus we only list them stating theirs relevant properties.

Let $(M,g)$ be an oriented $n$-dimensional Riemannian manifold and let $G\subset {\rm SO}(n)$ be the closed subgroup of the form
\begin{gather*}
G={\rm SO}(m)\times {\rm SO}(n-m)
\end{gather*}
for some $m=1,\ldots,n-1$. The quotient ${\rm SO}(n)/G$ is a symmetric space, which is the oriented Grassmannian $\operatorname{Gr}^{\rm o}_m(\mathbb{R}^n)$ of oriented $m$-dimensional subspaces in the Euclidean space $\mathbb{R}^n$. The reduction of the oriented orthonormal frame bundle ${\rm SO}(M)$ to the subbundle $P$ with the structure group~$G$ is equivalent to the existence of $m$-dimensional distribution $E$ on $M$ and, hence, its orthogonal complement $F=E^{\bot}$. We call $M$ with the distinguished distribution $E$ an almost product structure.

The connection $\nabla'$ induced by the connection form $\omega_{\mathfrak{g}}$, where $\omega$ is the connection form of the Levi-Civita connection $\nabla$, takes the form
\begin{gather*}
\nabla '_XY=\big(\nabla_XY^{\top}\big)^{\top}+\big(\nabla_XY^{\bot}\big)^{\bot},
\end{gather*}
where the decomposition $X=X^{\top}+X^{\bot}$ is taken with respect to $TM=E\oplus F$. In other words it is the sum of two connections induced by $\nabla$-connections in the vector bundles $E$ and $F$ over~$M$. The intrinsic torsion of almost product structure equals
\begin{gather*}
\xi_XY=\big(\nabla_XY^{\bot}\big)^{\top}+\big(\nabla_XY^{\top}\big)^{\bot}.
\end{gather*}
The associated bundle $N={\rm SO}(M)/G={\rm SO}(M)\times_{{\rm SO}(n)}({\rm SO}(n)/G)$ is the Grassmann bundle $\operatorname{Gr}^{\rm o}_m(TM)$ of $m$-dimensional oriented subspaces of tangent spaces to $M$ and the induced section $\sigma\colon M\to N$ is just the Gauss map of the distribution $E$. Therefore the main result (Theorem~\ref{thm:minimalityofP}) states that a ${\rm SO}(m)\times {\rm SO}(n-m)$-structure $P\subset {\rm SO}(M)$ is minimal if and only if the Gauss map $\sigma\colon (M,\tilde{g})\to (\operatorname{Gr}^{\rm o}_m(TM),\skal{\cdot}{\cdot})$ of $E$ is a harmonic map or, alternatively, if the image $E=\sigma(M)\subset \operatorname{Gr}^{\rm o}_m(TM)$ is minimal in the oriented Grassmann bundle $\operatorname{Gr}^{\rm o}_m(TM)$.

Let us now give two examples (compare~\cite{ggv,gd,hl}).

\begin{Example}
Let $(M,g,X_0)$ be a $K$-contact manifold with the Reeb vector f\/ield $X_0$, i.e., $X_0$~is a~unit Killing vector f\/ield and there exist one-form $\eta$ and endomorphism $\varphi$ such that
\begin{gather*}
\eta(X)=g(X,X_0),\qquad \varphi^2X=-X+\eta(X)X_0,\qquad d\eta(X,Y)=g(X,\varphi Y),\qquad \iota_{X_0}d\eta=0
\end{gather*}
for all $X,Y\in\Gamma(TM)$.

The one-dimensional distribution $E$ tangent to $X_0$ def\/ines the almost product structure on~$M$. Since $X_0$ is geodesic vector f\/ield, the distribution $E$ is totally geodesic. It can be shown, that $E\subset \operatorname{Gr}^{\rm o}_{n-1}(TM)$ is a minimal submanifold or, in other words, the immersion $\sigma\colon (M,\tilde{g})\to (\operatorname{Gr}^{\rm o}_{n-1}(TM),\skal{\cdot}{\cdot})$, $\sigma(x)=E_x$, $x\in M$, is minimal. Therefore, the ${\rm SO}(n-1)\times {\rm SO}(1)$-structure~$P$ induced by~$E$ is minimal in the orthonormal frame bundle~${\rm SO}(M)$.
\end{Example}

\begin{Example}
Consider the sphere $S^{4n-1}$ in the Euclidean space $\mathbb{R}^{4n}$ and let $I$, $J$, $K$ be the usual quaternionic structure on $\mathbb{R}^{4n}$. The $3$-dimensional subspaces spanned by $IN$, $JN$, $KN$, where $N$ is the unit outward vector f\/ield to $S^{4n-1}$, determine the Hopf distribution $E$ on $S^{4n-1}$. Since the Hopf distribution def\/ines the minimal immersion of $S^{4n-1}$ into the Grassmann bundle $\operatorname{Gr}^{\rm o}_3(TS^{4n-1})$, it follows that the ${\rm SO}(3)\times {\rm SO}(4n-4)$-structure $P\subset {\rm SO}(S^{4n-1})$ is minimal.
\end{Example}

\vspace{-2mm}

\pdfbookmark[1]{References}{ref}

\begin{thebibliography}{99}
\footnotesize\itemsep=0pt

\bibitem{lb}
Besse A.L., Einstein manifolds, \href{http://dx.doi.org/10.1007/978-3-540-74311-8}{\textit{Classics in Mathematics}}, Springer-Verlag,
  Berlin, 2008.

\bibitem{cg}
Chinea D., Gonzalez C., A classif\/ication of almost contact metric manifolds,
  \href{http://dx.doi.org/10.1007/BF01766972}{\textit{Ann. Mat. Pura Appl.~(4)}} \textbf{156} (1990), 15--36.

\bibitem{chs}
Chiossi S., Salamon S., The intrinsic torsion of {$\rm SU(3)$} and {${\rm
  G}_2$} structures, in Dif\/ferential Geometry ({V}alencia, 2001), \href{http://dx.doi.org/10.1007/BF01766972}{World Sci.
  Publ.}, River Edge, NJ, 2002, 115--133, \href{http://arxiv.org/abs/math.DG/0202282}{math.DG/0202282}.

\bibitem{cs}
Cleyton R., Swann A., Einstein metrics via intrinsic or parallel torsion,
  \href{http://dx.doi.org/10.1007/s00209-003-0616-x}{\textit{Math.~Z.}} \textbf{247} (2004), 513--528, \href{http://arxiv.org/abs/math.DG/0211446}{math.DG/0211446}.

\bibitem{dm}
Davidov J., Mu{\v{s}}karov O., Harmonic almost-complex structures on twistor
  spaces, \href{http://dx.doi.org/10.1007/BF02785864}{\textit{Israel~J. Math.}} \textbf{131} (2002), 319--332.

\bibitem{el}
Eells Jr. J., Sampson J.H., Harmonic mappings of {R}iemannian manifolds,
  \href{http://dx.doi.org/10.2307/2373037}{\textit{Amer.~J. Math.}} \textbf{86} (1964), 109--160.

\bibitem{ffs}
Falcitelli M., Farinola A., Salamon S., Almost-{H}ermitian geometry,
  \href{http://dx.doi.org/10.1016/0926-2245(94)00016-6}{\textit{Differential Geom. Appl.}} \textbf{4} (1994), 259--282.

\bibitem{fg}
Fern{\'a}ndez M., Gray A., Riemannian manifolds with structure group {${\rm
  G}_{2}$}, \href{http://dx.doi.org/10.1007/BF01760975}{\textit{Ann. Mat. Pura Appl.~(4)}} \textbf{132} (1982), 19--45.

\bibitem{gm}
Gil-Medrano O., Relationship between volume and energy of vector f\/ields,
  \href{http://dx.doi.org/10.1016/S0926-2245(01)00053-5}{\textit{Differential Geom. Appl.}} \textbf{15} (2001), 137--152.

\bibitem{ggv}
Gil-Medrano O., Gonz{\'a}lez-D{\'a}vila J.C., Vanhecke L., Harmonicity and
  minimality of oriented distributions, \href{http://dx.doi.org/10.1007/BF02803502}{\textit{Israel~J. Math.}} \textbf{143}
  (2004), 253--279.

\bibitem{gd}
Gonz{\'a}lez-D{\'a}vila J.C., Harmonicity and minimality of distributions on
  {R}iemannian manifolds via the intrinsic torsion, \href{http://dx.doi.org/10.4171/RMI/777}{\textit{Rev. Mat. Iberoam.}}
  \textbf{30} (2014), 247--275.

\bibitem{gmc2}
Gonz{\'a}lez-D{\'a}vila J.C., Mart{\'{\i}}n~Cabrera F., Harmonic
  {$G$}-structures, \href{http://dx.doi.org/10.1017/S0305004108001709}{\textit{Math. Proc. Cambridge Philos. Soc.}} \textbf{146}
  (2009), 435--459.

\bibitem{gh}
Gray A., Hervella L.M., The sixteen classes of almost {H}ermitian manifolds and
  their linear invariants, \href{http://dx.doi.org/10.1007/BF01796539}{\textit{Ann. Mat. Pura Appl.~(4)}} \textbf{123}
  (1980), 35--58.

\bibitem{hl}
Han D.-S., Lee E.-H., Harmonic {G}auss map and {H}opf f\/ibrations,
  \textit{J.~Korea Soc. Math. Educ. Ser.~B Pure Appl. Math.} \textbf{5} (1998),
  55--63.

\bibitem{mc}
Mart{\'{\i}}n~Cabrera F., Special almost {H}ermitian geometry, \href{http://dx.doi.org/10.1016/j.geomphys.2005.01.004}{\textit{J.~Geom.
  Phys.}} \textbf{55} (2005), 450--470, \href{http://arxiv.org/abs/math.DG/0409167}{math.DG/0409167}.

\bibitem{ms}
Mart{\'{\i}}n~Cabrera F., Swann A., The intrinsic torsion of almost
  quaternion-{H}ermitian manifolds, \href{http://dx.doi.org/10.5802/aif.2390}{\textit{Ann. Inst. Fourier (Grenoble)}}
  \textbf{58} (2008), 1455--1497, \href{http://arxiv.org/abs/0707.0939}{arXiv:0707.0939}.

\bibitem{an}
Naveira A.M., A classif\/ication of {R}iemannian almost-product manifolds,
  \textit{Rend. Mat.~(7)} \textbf{3} (1983), 577--592.

\bibitem{kn}
Niedzia{\l}omski K., On the frame bundle adapted to a submanifold,
  \href{http://dx.doi.org/10.1002/mana.201400041}{\textit{Math. Nachr.}} \textbf{288} (2015), 648--664, \href{http://arxiv.org/abs/1311.6172}{arXiv:1311.6172}.

\bibitem{vw}
Vergara-Diaz E., Wood C.M., Harmonic almost contact structures, \href{http://dx.doi.org/10.1007/s10711-006-9112-x}{\textit{Geom.
  Dedicata}} \textbf{123} (2006), 131--151, \href{http://arxiv.org/abs/math.DG/0602533}{math.DG/0602533}.

\bibitem{cw3}
Wood C.M., Harmonic almost-complex structures, \textit{Compositio Math.}
  \textbf{99} (1995), 183--212.

\bibitem{cw2}
Wood C.M., Harmonic sections of homogeneous f\/ibre bundles, \href{http://dx.doi.org/10.1016/S0926-2245(03)00021-4}{\textit{Differential
  Geom. Appl.}} \textbf{19} (2003), 193--210.

\end{thebibliography}
\end{document}
